%! Author = lili
%! Date = 7/13/26

\chapter{Tutorium 13.07.2026}\label{ch:tutorium-13.07.2026}


\section{PÜ 2.6.4}\label{sec:pu-2.6.4}
Ein Politiker fällt wiederholt durch fragwürdige Äußerungen auf.
Während er bei der letzten
Wahl noch 30\% der Stimmen erhalten hat, geht sein Wahlkampfteam mittlerweile davon aus,
daß dieser Wert gesunken ist.
Der Politiker widerspricht dem vehement.

In einer Stichprobe werden 200 Personen befragt.
45 Personen sprechen sich weiterhin für den
Politiker aus.

\begin{auf}
    \subsubsection*{Läßt sich die Hypothese, daß sich die Unterstützung der Wähler für den Politiker nicht verändert
    hat, halten? Verwenden Sie einen geeigneten Hypothesentest zum Signifikanzniveau von 5\%.}
    \paragraph{Parameter (Tutor):}
    $H_0 = \theta \in \Theta_0 = \left\{  \frac{3}{10} \right\}$ und
    $H_A = \theta \in \Theta_A = \left[ 0, \frac{3}{10}  \right)$.
    Man muss dann einen Hypothesen\gls{Hypothesentest} mit \gls{Signifikanzniveau} $\alpha = 0,05$.
    \paragraph{Modell dahinter (extra Info für Kontext, noch von der Aufgabe gefordert):}
    \begin{itemize}
        \item Stichprobenraum: $\chi = \{0,1\}^n$ mit $n = 200$
        \begin{itemize}
            \item Jedes $x in \chi$ ist ein n-fach unabhängig wiederholtes Bernoulli-Experiment
        \end{itemize}
        \item $\Theta = \left[ 0,1 \right] \leftarrow$ ``Zustimmung''
        \item $P = \{ \pe_{\theta} : \theta \in \Theta \}$, wobei
        \begin{itemize}
            \item $\pe_{\theta}$ durch die \gls{wfk} $\wfk_{\theta}$ mit $\wfk_{\theta}(x) = \theta^{\sumEN x_i} (1-\theta)^{n - \sumEN x_i}, \quad \forall x \in \chi$ gegeben ist.
        \end{itemize}
        \item Konkrete Stichprobe: $\tilde x \in \chi $ mit $\sumEN \tilde x_i = 45$
        \item $n = 200$
        \item Stichprobenvariablen:
        \begin{itemize}
            \item $X = (X_1, \dots, X_n)$ mit $\{X_i\}^n_{i=1}$ unabhängig identisch verteilt
            \item $X_i \sim Ber()$
        \end{itemize}
    \end{itemize}
    \paragraph{Weitere Lösung:}
    % Teststatistik
    \[
        t : \chi \rightarrow \{0, \dots, n\}
    \]
    Möglich:
    \[
        t(x) = n \ol{x} = \sumEN x_i
    \]
    \[
        t(x) = n - n \ol{x} = n - \sumEN x_i
    \]
    ABER: $ t(x) = n \ol{x} = \sumEN x_i$ dann wöre das bei Zustimmung irgendwie sehr groß, also nehmen wir:
    \[
        t(x) = n - n \ol{x} = n - \sumEN x_i
    \]
    deshalb
    % p-Wert & Test + E
    \begin{align*}
        p(x) & = sup_{\theta \in \Theta_0} \pe_{\theta}(t(X) \geq t(x)) \\
        % ja, erst großen, dann kleines x
        & = \pe_{\frac{3}{10}} (t(X) \geq t(x)) \\
        & = \pe_{\frac{3}{10}} (n - \sumEN X_i \geq n - \sumEN x_i) \\
        & = \pe_{\frac{3}{10}} (\sumEN X_i \leq \underbrace{ \sumEN x_i}_{n \ol{x}}) \\
    \end{align*}
    $\phi(x) = 1 \Rightarrow H_0$ wird verworfen
    \txc{(Die konkrete Stichprobe soll 1 sein, wenn wir unsere Stichprobe verwerfen. )}
    \[
        \phi(x) = \begin{cases}
                      1, & p(x) \leq \alpha \\ % nicht signifikant
                      0, & sonst
        \end{cases} = \mathbb{1}_{[0, \alpha]}
    \]
    % Durchführung
    Betrachte $p(\tilde x) = \pe_{\frac{3}{10}} (\sumEN X_i \leq 45)$
    % Erfolgsparameter 3/10, man geht davon aus alle dier Vernoulli Experimente haben den gleichen Erfolgsparameter
    \begin{align*}
        \sumEN X_i \sim Bin(n, \frac{3}{10}) & = Bin_{200, \frac{3}{10}}(\{0, \dots, 45\}) \\
        & = \sum^{45}_{i = 1} \binom{n}{1} \left( \frac{3}{10} \right)^i \left( \frac{7}{10} \right)^{n-i} \\
        & \approx 0,011 & Taschenrechner
    \end{align*}
    \[
        \Rightarrow \phi(\tilde x) % konkrete stichprobe
        = \mathbb{1}_{[0, 0.05]} (p(\tilde x)) % Indikatorfunktion
        = \mathbb{1}_{[0, \alpha]} (0,011) = 1
    \]
    % Schlussfolgerung
    $H_0$ muss verworfen werden.
\end{auf}

Ablauf so einer Aufgabenlösung:
\begin{enumerate}
    \item Hypothese aufstellen
    \item Signifikanzniveau festlegen
    \item Teststatistik aufstellen
    \item  p-Wert und Test + Entscheidungsregel
    \item ``Durchführung'' / Auswertung der Stichprobe
    \item Schlussfolgerung
\end{enumerate}

Hier umgesetzt als:
\begin{enumerate}
    \item Hypothese aufstellen
    \begin{itemize}
        \item $H_0 = \theta \in \Theta_0 = \left\{  \frac{3}{10} \right\}$ (Punkthypothese, keine Veränderung)
        \item $H_A = \theta \in \Theta_A = \left[ 0, \frac{3}{10}  \right)$ (Alternative Hypothese, negative
        Veränderung)
    \end{itemize}
    \item Signifikanzniveau festlegen
    \item Teststatistik aufstellen
    \item p-Wert und Test + Entscheidungsregel
    \item ``Durchführung'' / Auswertung der Stichprobe
    \item Schlussfolgerung
\end{enumerate}

Zusatz des Tutors:
\begin{mytable}
    \begin{tabular}{cccccc}
        k = & 45 & 46 & 47 & 48 & 49 \\
        $\pe_{\frac{3}{10}} (\sumEN X_i \leq k)$ & 0,011 & 0,017 & 0,025 & 0,036 & 0,056 \\
    \end{tabular}
    \caption{Ab 49 Zustimmungen hätte man die Hypothese nicht verwefen müssen - ``48 und kleiner'' ist unser Verwerfungsbereich}
\end{mytable}

% TODO konkrete Stickprobe vs. Stichproben
% TODO Arten von Hypothesen
\section{PÜ 2.6.5}\label{sec:pu-2.6.5}
Beantworten Sie die folgenden Fragen zur Theorie der \gls{Hypothesentest}s. Eine Begründung Ihrer
Antworten ist nicht nötig.
\begin{auf}
    \subsubsection*{Ein Fehler erster Art tritt ein, wenn ...}
    \paragraph{Lösung:} wir H0 verwerfen, wenn H0 wahr ist (Definition \gls{Fehler erster Art})
\end{auf}
Anmerkung: Fehler zweiter Art ist wir verwerfen H0 nicht, obwohl H0 falsch ist (wir H0 akzeptieren, wenn H0 falsch ist)

\komm{Es KANN passieren, dass Definitionen abgefragt werden in der Klausur (in Form von Multiple Choice oder speziellen Fragen). Oder man einen Hypothesentest aufstellen muss.}

\begin{auf}
    \subsubsection*{Wenn eine Hypothese zum Signifikanzniveau 0.6 verworfen wird, dann ...}
    \paragraph{Herleitung:}
    Entspricht ja
   \[
       \phi(x) = \begin{cases}
                     1, & p(x) \leq 0.6 \\ % nicht signifikant
                     0, & sonst
       \end{cases} = \mathbb{1}_{[0, 0.6]}
   \]
    Muss nicht für jedes Alpha verworfen werden, daher auch nicht für 0,5.
    \paragraph{Lösung:} kann es der Fall sein, dass sie zum Signifikanzniveau 0.5 verworfen wird
\end{auf}

\begin{auf}
    \subsubsection*{Der p-Wert eines Hypothesentests gibt folgendes an:}
    \paragraph{Lösung:} Das kleinste Signifikanzniveau für das Verwerfen der Hypothese H0 (per Definition, denn wenn
    wir eine konkrete Stichprobe haben (also p(x festhalten)) und das Signifikanzniveau danach aussuchen, dann kann
    das SN nicht kleiner als p(x) sein?)
\end{auf}

% evtl. kann man Statistik auf Lücke lernen...

\begin{auf}
    \subsubsection*{Ordnen Sie die folgenden Schritte gemäß der korrekten Durchführung eines Hypothesentests}
    \paragraph{Lösung:}
    \begin{enumerate}
        \item Wahl von Hypothese und Alternative
        \item Wahl des Signifikanzniveaus
        \item Wahl der Teststatistik
        \item Wahl der Entscheidungsregel
        \item Auswertung der Stichprobe
        \item Schlussfolgerung
    \end{enumerate}
\end{auf}

\section{Besprechung ÜB 2.5.1}\label{sec:besprechung-ub-2.5.1}
Wir nehmen in dieser Aufgabe an, daß unsere Stichprobe von einem n-mal unabhängig wieder-
holten Bernoulli-Experiment mit unbekannter Erfolgswahrscheinlichkeit $\theta \in (0, 1)$ stammt.

Wir wissen aus der Vorlesung, daß das arithmetische Mittel der Stichprobenwerte ein erwartungstreuer Schätzer für $\theta$ ist.

\[
todo
\]

\begin{auf}
    \subsubsection*{Begründen Sie, warum t so definiert ein Schätzer für die Varianz $\theta(1- \theta)$ ist.}
    Siehe Bild TODO % TODO Bild
    \subsubsection*{Ist dieser Schätzer erwartungstreu?}
\end{auf}
